\chapter{归约}
\label{chap:reduction}

\setcounter{footnote}{0}

归约（reduction）操作从一列值中导出一个值。这个结果可以是所有输入值的和、最大值、
最小值等，输入与结果也可以采用整数、单精度浮点数、双精度浮点数、半精度浮点数、
字符等不同类型。各种归约都有相同的计算结构。与直方图一样，归约也是一种重要的
计算模式，因为它能从大量数据中提炼出概括信息。并行归约要求并行线程相互协调，
才能得到正确结果。协调方式必须精心设计，否则很容易触发并行计算系统中常见的性能
瓶颈。因此，并行归约很适合用来说明这些瓶颈，并引出缓解它们的技术。

\section{背景}
\label{sec:reduction-background}

归约根据某个二元运算符，把一组元素变换为单个元素。例如，可以用加法运算符定义
浮点数集合上的求和归约，结果就是集合中所有浮点数之和。对集合
$\{7.0,2.1,5.3,9.0,11.2\}$，求和归约得到

\[
  7.0+2.1+5.3+9.0+11.2=34.6.
\]

归约可以顺序遍历数组中的每个元素来完成。图 \ref{fig:10-1} 给出了一个简单的 C
语言顺序求和归约。代码先用加法的单位元 $0.0$ 初始化结果变量 \code{sum}，
\footnote{二元运算符的单位元作为一个操作数时，结果总是另一个操作数。例如，浮点
加法的单位元为 $0.0$，因为任意浮点值 $v$ 与 $0.0$ 相加仍为 $v$。代码中使用单位元
只是为了方便；如果运算符没有单位元，可以用元素 0 初始化累加量，并让循环从
\code{i=1} 开始。}再用 for 循环遍历保存输入值的数组。在第 $i$ 次迭代中，代码把
\code{sum} 的当前值与输入数组的第 $i$ 个元素相加。在本例中，迭代 0 后，
\code{sum} 为 $0.0+7.0=7.0$；迭代 1 后，它为 $7.0+2.1=9.1$。也就是说，每次
迭代都把一个输入元素累加到 \code{sum} 中。处理完 5 个元素后，\code{sum} 为
$34.6$，即归约结果。

\begin{figure}[htbp]
  \centering
  \begin{minipage}{0.88\textwidth}
  \begin{lstlisting}[language=C++,numbers=left]
sum = 0.0f;
for(i = 0; i < N; ++i) {
    sum += input[i];
}
  \end{lstlisting}
  \end{minipage}
  \caption{简单的顺序求和归约代码}
  \label{fig:10-1}
\end{figure}

\noindent\textbf{译者注：}底本在这个 5 元素示例中写的是“迭代 5 后”，但循环下标
从 0 开始，实际最后一次迭代是迭代 4。本译稿改为“处理完 5 个元素后”，以免把
元素个数与迭代下标混淆。

还可以为许多其他运算符定义归约。以单位元 $1.0$ 初始化浮点乘法运算符，就得到乘积
归约，其结果是集合中所有浮点数的乘积。最小值归约可以用返回两个输入中较小者的
比较运算符定义；最大值归约则使用返回较大者的比较运算符。对实数而言，最小值运算符
的单位元是 $+\infty$，最大值运算符的单位元是 $-\infty$。因此，浮点数集合的最小值
（最大值）归约会生成列表中的最小（最大）数。

\noindent\textbf{译者注：}底本把上段写成“minimum (minimum) operator”，但前后文
明确分别讨论最小值与最大值运算符。本译稿按上下文改为“最小值（最大值）运算符”。

图 \ref{fig:10-2} 给出了任意运算符上的顺序归约通式。这里的运算符定义为一个接收
两个输入并返回一个值的函数。循环访问某个元素时，执行的动作取决于归约类型。例如，
最大值归约的 \code{Operator} 函数比较两个输入并返回较大者；最小值归约则返回较小者。
for 循环访问完所有元素后，顺序算法结束。对包含 $N$ 个元素的集合，循环执行 $N$ 次，
退出时得到归约结果。

\begin{figure}[htbp]
  \centering
  \begin{minipage}{0.88\textwidth}
  \begin{lstlisting}[language=C++,numbers=left]
acc = IDENTITY;
for(i = 0; i < N; ++i) {
    acc = Operator(acc, input[i]);
}
  \end{lstlisting}
  \end{minipage}
  \caption{顺序归约代码的通式}
  \label{fig:10-2}
\end{figure}

归约非常常用，一些编程语言已经把它纳入标准库。例如，C++ 标准库中的
\code{std::accumulate} 和 \code{std::reduce} 可以执行不同形式的归约。类似地，
Thrust\cite{bell2012thrust,nvidia-thrust2025} 和 CUB\cite{nvidia-cub2025} 等 CUDA
库也提供了针对 GPU 高度优化的并行归约实现。尽管如此，本章仍从头讲解如何在 GPU
上并行化并优化归约，因为它能清楚展示许多并行化和优化技术。

\section{归约树}
\label{sec:reduction-trees}

并行归约算法在文献中已有大量研究。图 \ref{fig:10-3} 展示了它的基本思想：时间沿
竖直方向向下推进，每个时间步中线程并行执行的活动沿水平方向排列。

\begin{figure}[htbp]
  \centering
  \scriptsize
  \begin{tabular}{c@{\hspace{3em}}c}
    \begin{tabular}{c}
      \fbox{3}\ \fbox{1}\ \fbox{7}\ \fbox{0}\ \fbox{4}\ \fbox{1}\ \fbox{6}\ \fbox{3}\\[0.35em]
      $\searrow\!\max\!\swarrow\quad\searrow\!\max\!\swarrow\quad
       \searrow\!\max\!\swarrow\quad\searrow\!\max\!\swarrow$\\[-0.15em]
      \fbox{3}\qquad\fbox{7}\qquad\fbox{4}\qquad\fbox{6}\\[0.35em]
      $\searrow\!\max\!\swarrow\qquad\qquad\searrow\!\max\!\swarrow$\\[-0.15em]
      \fbox{7}\hspace{4.8em}\fbox{6}\\[0.35em]
      $\searrow\!\max\!\swarrow$\\[-0.15em]
      \fbox{7}\\[0.4em]
      (a) 并行最大值归约树
    \end{tabular}
    &
    \begin{tabular}{c}
      \fbox{4}\ \fbox{7}\ \fbox{2}\ \fbox{3}\ \fbox{8}\ \fbox{5}\ \fbox{9}\ \fbox{6}\\[0.35em]
      $\searrow\!+\!\swarrow\quad\searrow\!+\!\swarrow\quad
       \searrow\!+\!\swarrow\quad\searrow\!+\!\swarrow$\\[-0.15em]
      \fbox{11}\qquad\fbox{5}\qquad\fbox{13}\qquad\fbox{15}\\[0.35em]
      $\searrow\!+\!\swarrow\qquad\qquad\searrow\!+\!\swarrow$\\[-0.15em]
      \fbox{16}\hspace{4.3em}\fbox{28}\\[0.35em]
      $\searrow\!+\!\swarrow$\\[-0.15em]
      \fbox{44}\\[0.4em]
      (b) 并行求和归约树
    \end{tabular}
  \end{tabular}
  \caption{最大值归约与求和归约的并行归约树示例}
  \label{fig:10-3}
\end{figure}

先看图 \ref{fig:10-3}(a) 的最大值归约树。第一轮（时间步）并行处理原始元素的四对
组合，执行四次最大值运算，得到四个部分归约结果，也就是每一对中的较大值。第二个
时间步并行处理这四个部分结果的两对组合，执行两次最大值运算，生成两个更接近最终
结果的部分结果。它们分别是原输入前四个元素和后四个元素中的最大值。第三个也是最后
一个时间步执行一次最大值运算，得到原输入中的最大值 7。

从顺序归约改为并行归约时，运算的执行次序发生了变化。以图 \ref{fig:10-3}(a) 顶部
的输入为例，图 \ref{fig:10-2} 那样的顺序最大值归约会先比较 \code{acc} 中的单位元
$-\infty$ 与输入值 3，把较大的 3 写回 \code{acc}；再比较 \code{acc} 中的 3 与
输入值 1，结果仍为 3；随后比较 3 与 7，得到 7。图 \ref{fig:10-3}(a) 的并行归约则
先比较输入值 7 与输入值 0，然后才把结果与 3 和 1 的最大值比较。要并行化归约，
关键就在于能够这样改变运算次序。

为了重排运算，必须假定把运算符应用于输入值的顺序不会影响结果。最大值归约无论以
什么顺序应用运算符，结果都相同。运算符满足\textbf{结合律（associativity）}时，
数学上可以保证这一性质。若运算符 $\Theta$ 满足

\[
  (a\mathbin{\Theta}b)\mathbin{\Theta}c
  =a\mathbin{\Theta}(b\mathbin{\Theta}c),
\]

就称它具有结合性。换句话说，对于只含该运算符的表达式，可以在任意位置加括号而
不改变结果。借助这种等价关系，我们可以在保持结果等价的同时，把一种运算次序变成
另一种。

例如，整数加法满足结合律，即 $(1+2)+3=1+(2+3)$。\footnote{从数学意义上说，
整数加法满足结合律；但在 C++ 中并非绝对如此，因为运算次序可能影响是否发生溢出或
下溢，而有符号整数溢出的行为未定义。如果程序员确定不会溢出或下溢，实际编程时仍可
把整数加法视为满足结合律。}整数减法则不满足结合律，因为
$(1-2)-3\ne 1-(2-3)$。可以把减法变为与第二个操作数的相反数相加，从而获得结合性；
例如，$(1-2)-3$ 可以改写为
$(1+(-2))+(-3)=1+((-2)+(-3))$。

对实数而言，加法在数学意义上满足结合律；但严格来说，由于不同加括号方式可能产生
不同的舍入结果，浮点加法并不满足结合律。有些应用只要不同浮点结果之间的差异不超过
容许范围，就认为结果相同。凭借这种容差，开发者和编译器优化在实践中可以把浮点加法
视为具有结合性。附录 A 将详细讨论这个问题。

把图 \ref{fig:10-2} 的顺序归约转换为图 \ref{fig:10-3}(a) 的归约树，要求运算符满足
结合律。可以把归约看成一列运算；二者的次序差异来自在同一运算列表的不同位置加入
括号。图 \ref{fig:10-2} 对应

\[
  ((((((3\mathbin{\max}1)\mathbin{\max}7)\mathbin{\max}0)
  \mathbin{\max}4)\mathbin{\max}1)\mathbin{\max}6)\mathbin{\max}3,
\]

而图 \ref{fig:10-3}(a) 对应

\[
  ((3\mathbin{\max}1)\mathbin{\max}(7\mathbin{\max}0))
  \mathbin{\max}((4\mathbin{\max}1)\mathbin{\max}(6\mathbin{\max}3)).
\]

第 \ref{sec:reduction-control-divergence} 节还会采用一种优化，它不仅调整应用运算符的
次序，还会调整操作数的次序。要重排操作数，运算符还必须满足\textbf{交换律
（commutativity）}。若 $a\mathbin{\Theta}b=b\mathbin{\Theta}a$，就称运算符
$\Theta$ 具有交换性；也就是说，调换表达式中操作数的位置仍得到相同结果。最大值
运算符以及最小值、加法和乘法等许多运算符都满足交换律。\footnote{对实数而言，最大值
和最小值在数学意义上满足交换律；但对浮点数并非绝对如此，因为其中一个操作数为 NaN
时会有特殊行为。详细规则见 IEEE 浮点标准。如果程序员确定输入中没有 NaN，实际编程
时可以把最大值和最小值视为满足交换律。}显然，并非所有运算符都如此。例如，加法满足
$1+2=2+1$，整数减法却不满足 $1-2=2-1$。

图 \ref{fig:10-3}(a) 的并行模式称为\textbf{归约树（reduction tree）}，因为它看起来
像一棵以原始输入元素为叶、以最终结果为根的树。这里的归约树不要与树形数据结构混淆；
后者通过指针或约定的位置显式或隐式连接结点，归约树中的边则只是概念关系，表示信息
从一个时间步中的运算流向下一个时间步。

为了比较顺序归约与并行归约实现的速度，需要引入\textbf{工作量（work）}和
\textbf{跨度（span）}。实现的工作量是它执行的运算数；算法的跨度则是在执行资源无限
时完成这些运算所需的时间。例如，第 2 章的向量加法内核为两个长度为 $N$ 的向量启动
$N$ 个线程。执行时，每个线程完成一次操作，包括读取两个输入元素、执行加法并写出一个
输出元素。因此，该内核完成 $O(N)$ 的工作量，即所有线程合计执行 $N$ 次操作。

把一个线程完成一次操作的时间称为一个\textbf{步（step）}。涉及多个线程时，一个步
是所有线程或其中一部分线程各自并行执行一次操作的时间段。计算的跨度用实现所需的步数
衡量，并假定执行资源数量不受限制。顺序向量加法的跨度为 $O(N)$，因为 $N$ 次操作依次
执行；并行向量加法内核的跨度则为 $O(1)$，因为 $N$ 次操作全部并行执行。

分析顺序算法时，工作量与跨度没有区别，因为操作依次执行，算法所需步数与工作量直接
成正比。例如，顺序向量加法完成 $O(N)$ 的工作，也需要 $O(N)$ 个步。

顺序归约同样完成 $O(N)$ 的工作并需要 $O(N)$ 个步。相比之下，图
\ref{fig:10-3}(a) 的并行归约树用 7 次运算归约 8 个值：第一步 4 次，第二步 2 次，
第三步 1 次。一般来说，对 $N$ 个输入值，第一轮执行 $N/2$ 次，第二轮执行 $N/4$ 次，
依此类推。因此总运算数为等比级数

\[
  \frac{N}{2}+\frac{N}{4}+\frac{N}{8}+\cdots+\frac{N}{N}=N-1.
\]

换句话说，并行归约树的工作量仍为 $O(N)$，与顺序归约相同。第 11 章将看到，并行算法
并不总能保持与顺序算法相同的工作量，这会引出工作效率的概念。

跨度方面，归约 $N$ 个输入值的并行归约树需要 $\log_2N$ 个步。每个步执行一层运算，
所以跨度为 $O(\log N)$。例如，在执行资源充足时，只需 $\log_2 1024=10$ 个步就能归约
$N=1024$ 个输入值。第一步需要 $N/2=512$ 份执行资源，之后所需资源迅速减少，最后一步
只需要一份。每一步的并行度等于该步所需的执行单元数。所有步的平均并行度等于总运算数
除以步数，即 $(N-1)/\log_2N$。当 $N=1024$ 时，10 个步的平均并行度为 102.3，而峰值
并行度为第一步的 512。并行度与资源消耗随步骤大幅变化，使归约树成为并行计算系统中
颇具挑战性的模式。

\begin{figure}[htbp]
  \centering
  \scriptsize
  \begin{tabular}{c@{\hspace{1.2em}}c@{\hspace{1.2em}}c@{\hspace{1.2em}}c}
    \multicolumn{4}{c}{\textbf{四分之一决赛}}\\[0.2em]
    \fbox{荷兰 2:1 巴西} & \fbox{乌拉圭 1:1 加纳} &
    \fbox{阿根廷 0:4 德国} & \fbox{巴拉圭 0:1 西班牙}\\[-0.05em]
    7 月 2 日，纳尔逊·曼德拉湾 & 7 月 2 日，约翰内斯堡 &
    7 月 3 日，开普敦 & 7 月 3 日，约翰内斯堡\\[0.4em]
    \multicolumn{2}{c}{$\searrow$\quad\fbox{荷兰 3:2 乌拉圭}\quad$\swarrow$} &
    \multicolumn{2}{c}{$\searrow$\quad\fbox{德国 0:1 西班牙}\quad$\swarrow$}\\[-0.05em]
    \multicolumn{2}{c}{7 月 6 日，开普敦} & \multicolumn{2}{c}{7 月 7 日，德班}\\[0.15em]
    \multicolumn{4}{c}{\textbf{半决赛}}\\[0.4em]
    \multicolumn{4}{c}{$\searrow$\quad\fbox{荷兰 0:1 西班牙}\quad$\swarrow$}\\[-0.05em]
    \multicolumn{4}{c}{7 月 11 日，约翰内斯堡（加时）}\\[0.15em]
    \multicolumn{4}{c}{\textbf{决赛}\qquad$\Downarrow$\qquad\fbox{冠军：西班牙}}
  \end{tabular}
  \caption{作为归约树的 2010 年世界杯淘汰赛末段赛程}
  \label{fig:10-4}
\end{figure}

\paragraph{体育与竞赛中的并行归约。}
早在计算机出现之前，体育和竞赛就已经采用并行归约。图 \ref{fig:10-4} 展示了 2010
年南非世界杯四分之一决赛、半决赛和决赛的赛程，它实际上就是重新排布过的归约树。
世界杯的淘汰过程是一种最大值归约，其中“最大值”运算符返回击败另一支球队的球队。
赛事归约分多轮进行：球队两两配对，第一轮的各场比赛并行进行，胜者进入第二轮，
第二轮胜者再进入第三轮，以此类推。8 支球队参赛时，第一轮（四分之一决赛）产生
4 支胜队，第二轮（半决赛）产生 2 支胜队，第三轮（决赛）产生最终冠军。每一轮都是
归约过程中的一个步，因此图 \ref{fig:10-4} 的淘汰赛跨度为 3 个步。

\noindent\textbf{译者注：}底本写成“跨度为 4 个步”，但图中 8 支球队依次经过
四分之一决赛、半决赛和决赛三轮，且紧接着的正文也说明 1024 支球队需要
$\log_2 1024=10$ 轮。本译稿据此改为 3 个步。

即使有 1024 支球队，也只需 10 轮就能决出冠军。关键在于要有足够资源，让第一轮的
512 场比赛、第二轮的 256 场、第三轮的 128 场等分别并行进行。资源充足时，即使有
6 万支球队，也只需 16 轮就能产生最终胜者。归约树能大幅加速归约过程，但也消耗大量
资源。在世界杯的例子中，一场比赛需要大型足球场、裁判和工作人员，还需要酒店和餐馆
接待大量观众。图 \ref{fig:10-4} 的四场四分之一决赛在三个城市举行：纳尔逊·曼德拉湾
（伊丽莎白港）、开普敦和约翰内斯堡，三座城市合计提供了举办四场比赛所需的资源。
约翰内斯堡的两场比赛在不同日期进行，说明两场比赛共享资源会延长归约过程。计算中的
归约树也会遇到类似权衡。

图 \ref{fig:10-3}(b) 给出了 8 个输入值的求和归约树。前面对最大值归约树工作量与
跨度的分析同样适用于它。使用最多 4 个加法器及其读写数据所需资源，只需
$\log_2 8=3$ 个时间步即可完成归约。下面将用这个例子说明求和归约内核的设计。

\section{一个简单的归约内核}
\label{sec:simple-reduction-kernel}

现在可以着手实现图 \ref{fig:10-3}(b) 的并行求和归约树了。归约树需要所有线程协作，
而整个网格中的线程无法直接完成这种协作，所以先实现只在一个线程块内执行求和归约树
的内核。对包含 $N$ 个元素的输入数组，调用这个简单内核时启动一个网格，其中只有一个
线程块，块内有 $N/2$ 个线程。线程块最多可有 1024 个线程，因此这种做法最多能处理
2048 个输入元素。第 \ref{sec:reduction-arbitrary-length} 节将解除这个限制。

第一步中，全部 $N/2$ 个线程都参与计算，每个线程把两个元素相加，生成 $N/2$ 个部分
和。下一步中，一半线程退出，只剩 $N/4$ 个线程继续计算并生成 $N/4$ 个部分和。这个
过程不断重复，直到最后一步只剩一个线程，由它产生总和。

\begin{figure}[htbp]
  \centering
  \begin{minipage}{0.95\textwidth}
  \begin{lstlisting}[language=C++,numbers=left]
__global__ void SimpleSumReductionKernel(float* input, float* output) {
    unsigned int i = 2*threadIdx.x;
    for(unsigned int stride = 1; stride <= blockDim.x; stride *= 2) {
        if(threadIdx.x % stride == 0) {
            input[i] += input[i + stride];
        }
        __syncthreads();
    }
    if(threadIdx.x == 0) {
        *output = input[0];
    }
}
  \end{lstlisting}
  \end{minipage}
  \caption{简单的求和归约内核}
  \label{fig:10-5}
\end{figure}

图 \ref{fig:10-5} 给出了简单的求和内核，图 \ref{fig:10-6} 则展示了该代码实现的归约树
如何执行。图中时间从上向下推进。假设输入数组位于全局内存中，调用内核函数时传入
指向它的指针。每个线程分配到位置 \code{2*threadIdx.x}（第 2 行），也就是输入数组的
偶数位置：线程 0 对应 \code{input[0]}，线程 1 对应 \code{input[2]}，线程 2 对应
\code{input[4]}，依此类推。每个线程都是所分配位置的“所有者”，只有它可以写入该位置。
这种设计遵循所有者计算规则：每个数据位置只归一个线程所有，也只能由该线程更新。

\begin{figure}[htbp]
  \centering
  \scriptsize
  \setlength{\tabcolsep}{2.2pt}
  \begin{tabular}{c|*{16}{c}}
    &0&1&2&3&4&5&6&7&8&9&10&11&12&13&14&15\\\hline
    线程所有者&$T_0$&&$T_1$&&$T_2$&&$T_3$&&$T_4$&&$T_5$&&$T_6$&&$T_7$&\\
    步 0，$s=1$&\colorbox{blue!25}{\strut$\Sigma$}&&\colorbox{blue!25}{\strut$\Sigma$}&&
      \colorbox{blue!25}{\strut$\Sigma$}&&\colorbox{blue!25}{\strut$\Sigma$}&&
      \colorbox{blue!25}{\strut$\Sigma$}&&\colorbox{blue!25}{\strut$\Sigma$}&&
      \colorbox{blue!25}{\strut$\Sigma$}&&\colorbox{blue!25}{\strut$\Sigma$}&\\
    步 1，$s=2$&\colorbox{blue!25}{\strut$\Sigma$}&&&&
      \colorbox{blue!25}{\strut$\Sigma$}&&&&\colorbox{blue!25}{\strut$\Sigma$}&&&&
      \colorbox{blue!25}{\strut$\Sigma$}&&&\\
    步 2，$s=4$&\colorbox{blue!25}{\strut$\Sigma$}&&&&&&&&
      \colorbox{blue!25}{\strut$\Sigma$}&&&&&&&\\
    步 3，$s=8$&\colorbox{blue!25}{\strut$\Sigma$}&&&&&&&&&&&&&&&
  \end{tabular}
  \\[0.45em]
  \begin{tabular}{c@{\quad}l}
    步 0 & $\code{input[2k]}\leftarrow\code{input[2k]}+\code{input[2k+1]}$\\
    步 1 & $\code{input[4k]}\leftarrow\code{input[4k]}+\code{input[4k+2]}$\\
    步 2 & $\code{input[8k]}\leftarrow\code{input[8k]}+\code{input[8k+4]}$\\
    步 3 & $\code{input[0]}\leftarrow\code{input[0]}+\code{input[8]}$
  \end{tabular}
  \caption{线程“所有者”到输入数组位置的分配，以及简单归约内核随时间推进的执行过程；
  每一层对应一次循环迭代}
  \label{fig:10-6}
\end{figure}

图 \ref{fig:10-6} 顶行表示线程到输入数组位置的分配，后续每一行表示图
\ref{fig:10-5} 的 for 循环一次迭代结束后写入的位置。每次迭代中被覆盖的位置以
$\Sigma$ 标出。例如，第一次迭代结束时，偶数下标位置被原数组相邻元素对
（0--1、2--3、4--5 等）的部分和覆盖。第二次迭代结束时，下标为 4 的倍数的位置被
四个相邻原始元素（0--3、4--7 等）的部分和覆盖。

图 \ref{fig:10-5} 使用 \code{stride} 变量，让线程找到应当累加到所有者位置的部分和。
它从 1 开始（第 3 行），每次迭代翻倍，依次取 1、2、4、8 等，直到大于块中的线程总数
\code{blockDim.x}。图 \ref{fig:10-6} 中，每个活跃线程都把距离其所有者位置
\code{stride} 个元素的数组元素累加进去。例如，在迭代 0 中，线程 0 以步长 1 把
\code{input[1]} 加到自己的 \code{input[0]}，于是 \code{input[0]} 变成原始
\code{input[0]} 与 \code{input[1]} 的部分和。迭代 1 中，线程 0 以步长 2 把
\code{input[2]} 加到 \code{input[0]}。此时 \code{input[2]} 已保存原始
\code{input[2]} 和 \code{input[3]} 的和，\code{input[0]} 则保存原始
\code{input[0]} 和 \code{input[1]} 的和，因此迭代 1 结束后，\code{input[0]}
保存原数组前四个元素之和。最后一次迭代后，\code{input[0]} 包含全部原始元素之和，
也就是求和归约结果，线程 0 把它写入最终输出（第 10 行）。

\noindent\textbf{译者注：}底本把“块中的线程总数”写成 \code{blockIdx.x}，但该变量
表示块在网格中的下标；图 \ref{fig:10-5} 的循环条件实际使用 \code{blockDim.x}。
本译稿按代码改为后者。

再看图 \ref{fig:10-5} 第 4 行的 if 语句。它的条件用于选出每次迭代中的活跃线程。
图 \ref{fig:10-6} 表明，迭代 $n$ 中执行加法的是线程下标为 $2^n$ 的倍数的线程。
条件 \code{threadIdx.x \% stride == 0} 正是在检查线程下标是否为步长的倍数，而
\code{stride} 在各次迭代中依次为 $1,2,4,8,\ldots$，也就是迭代 $n$ 中的 $2^n$。
所有线程都执行相同的内核代码；满足条件的活跃线程执行第 5 行加法，不满足条件的非活跃
线程跳过它。迭代越往后，活跃线程越少；最后一次迭代只剩线程 0，由它生成归约结果。

循环中的 \code{__syncthreads()}（第 7 行）保证当前迭代计算出的所有部分和都已经写入
输入数组中的目标位置，任何线程才能开始下一次迭代。这样，进入下一次迭代的线程才能
正确读取上次迭代产生的部分和。例如，第一次迭代会用成对部分和替换偶数位置；
\code{__syncthreads()} 保证这些结果确实全部就绪，第二次迭代中的活跃线程才会使用它们。

还要注意，一个线程写入全局内存的数据不能在没有内存栅栏的情况下由其他线程读取。
这里跨迭代读写相同位置的线程都属于同一线程块，因此每次迭代末尾的
\code{__syncthreads()} 也充当块内线程的内存栅栏：调用它之前写入的数据，在调用完成后
可以被同一块中的线程正确读取。规定线程通过内存访问交换数据时应有何种可观察行为的
规则称为\textbf{内存一致性模型（memory consistency model）}。更多细节可参阅 CUDA
编程指南。

\section{减少控制流分歧}
\label{sec:reduction-control-divergence}

图 \ref{fig:10-5} 的代码实现了图 \ref{fig:10-6} 的并行归约树，也能产生预期的求和
结果。遗憾的是，它在各次迭代中管理活跃和非活跃线程的方式造成了严重的控制流分歧。
例如，图 \ref{fig:10-6} 中第二次迭代只有 \code{threadIdx.x} 为偶数的线程执行加法。
正如第 4 章所述，控制流分歧会显著降低执行资源利用效率，也就是用于生成有效结果的
资源比例。本例中，一个 warp 的 32 个线程都占用执行资源，实际却只有一半活跃，浪费了
一半资源。随着时间推进，浪费还会增加：第三次迭代只有四分之一线程活跃，浪费四分之三；
到第 6 次迭代，一个 warp 中只有 1 个线程活跃，浪费 $31/32$ 的执行资源。

如果输入数组长度大于 32，第 6 次迭代后会有整个 warp 变为非活跃。以 256 个输入元素
为例，内核启动 128 个线程，即 4 个 warp。前 6 次迭代中，四个 warp 都呈现上段所述的
分歧模式。第 7 次迭代时，warp 1 和 warp 3 完全非活跃，不再产生控制流分歧；warp 0
和 warp 2 则各只有一个活跃线程，浪费 $31/32$ 的执行资源。第 8 次迭代只有 warp 0
活跃，其中也只有一个线程工作。

一般来说，长度为 $N$ 的输入数组，其执行资源利用效率等于所有迭代中的活跃线程总数
与消耗的执行资源总量之比。后者正比于所有迭代中的活跃 warp 总数，因为只要 warp
仍活跃，无论内部有几个线程实际工作，它都要消耗完整的 warp 执行资源。总消耗为

\[
  \left(\frac{N}{64}\cdot 6+
  \frac{N}{64}\cdot\frac12+
  \frac{N}{64}\cdot\frac14+\cdots+1\right)\cdot 32.
\]

这里 $N/64$ 是启动的 warp 总数，因为内核启动 $N/2$ 个线程，每 32 个线程组成一个
warp。$N/64$ 乘以 6，是因为前 6 次迭代中所有 warp 都活跃；之后每次迭代活跃 warp
数减半。括号中的表达式就是跨全部迭代的活跃 warp 总数，外面的 32 表示每个活跃 warp
总会占用 32 个线程对应的完整执行资源。对 256 个输入元素，消耗量为
$(4\times6+2+1)\times32=864$。

活跃线程提交的有效结果数，就是所有迭代中的活跃线程总数：

\[
  \frac{N}{64}(32+16+8+4+2+1)
  +\frac{N}{64}\cdot\frac12\cdot1
  +\frac{N}{64}\cdot\frac14\cdot1+\cdots+1.
\]

括号中的各项给出前 6 次迭代里每个 warp 的活跃线程数。从第 7 次迭代开始，活跃 warp
数每次减半，每个活跃 warp 中只剩一个活跃线程。对 256 个输入元素，有效结果总数为
$4(32+16+8+4+2+1)+2+1=255$。这也很直观，因为把 256 个值归约为一个值恰好需要
255 次运算。

结合这两个结果，输入长度为 256 时，执行资源利用效率为 $255/864=0.30$。也就是说，
并行执行资源没有充分发挥加速潜力；平均只有约 30\% 的已消耗资源真正贡献于求和结果。

上述分析表明，控制流分歧跨越多个 warp，并且随时间持续存在。更好的线程与输入位置
分配方式可以减少分歧并提高利用效率。图 \ref{fig:10-6} 的问题在于，部分和的位置随
时间越来越分散，拥有这些位置的活跃线程也越来越分散，非活跃线程比例和资源闲置程度
因而不断上升。

\begin{figure}[htbp]
  \centering
  \scriptsize
  \setlength{\tabcolsep}{2.2pt}
  \begin{tabular}{c|*{16}{c}}
    &0&1&2&3&4&5&6&7&8&9&10&11&12&13&14&15\\\hline
    线程所有者&$T_0$&$T_1$&$T_2$&$T_3$&$T_4$&$T_5$&$T_6$&$T_7$&&&&&&&&\\
    步 0，$s=8$&\colorbox{blue!25}{\strut$\Sigma$}&\colorbox{blue!25}{\strut$\Sigma$}&
      \colorbox{blue!25}{\strut$\Sigma$}&\colorbox{blue!25}{\strut$\Sigma$}&
      \colorbox{blue!25}{\strut$\Sigma$}&\colorbox{blue!25}{\strut$\Sigma$}&
      \colorbox{blue!25}{\strut$\Sigma$}&\colorbox{blue!25}{\strut$\Sigma$}&&&&&&&&\\
    步 1，$s=4$&\colorbox{blue!25}{\strut$\Sigma$}&\colorbox{blue!25}{\strut$\Sigma$}&
      \colorbox{blue!25}{\strut$\Sigma$}&\colorbox{blue!25}{\strut$\Sigma$}&&&&&&&&&&&&\\
    步 2，$s=2$&\colorbox{blue!25}{\strut$\Sigma$}&\colorbox{blue!25}{\strut$\Sigma$}&&&&&&&&&&&&&&\\
    步 3，$s=1$&\colorbox{blue!25}{\strut$\Sigma$}&&&&&&&&&&&&&&&
  \end{tabular}
  \\[0.45em]
  \begin{tabular}{c@{\quad}l}
    步 0 & $\code{input[k]}\leftarrow\code{input[k]}+\code{input[k+8]}$，$0\le k<8$\\
    步 1 & $\code{input[k]}\leftarrow\code{input[k]}+\code{input[k+4]}$，$0\le k<4$\\
    步 2 & $\code{input[k]}\leftarrow\code{input[k]}+\code{input[k+2]}$，$0\le k<2$\\
    步 3 & $\code{input[0]}\leftarrow\code{input[0]}+\code{input[1]}$
  \end{tabular}
  \caption{减少控制流分歧的更优线程到输入数组位置分配方式}
  \label{fig:10-7}
\end{figure}

确实存在一种能显著减少控制流分歧的分配策略：安排线程及其所有者位置，使活跃线程随
时间推进仍彼此相邻。换言之，让步长随时间减小，而不是增大。图 \ref{fig:10-7} 对
16 个输入元素展示了新策略。线程分配到数组前半部分的位置。第一次迭代中，每个线程
跨过半个数组，把对应的另一个元素加到自己的位置：线程 0 把 \code{input[8]} 加到
\code{input[0]}，线程 1 把 \code{input[9]} 加到 \code{input[1]}，依此类推。

以后的每次迭代中，一半活跃线程退出，其余线程把距离所有者位置“当前活跃线程数”个
元素的输入加进来。例如迭代 3 时只剩线程 0 和线程 1：线程 0 把 \code{input[2]}
加到 \code{input[0]}，线程 1 把 \code{input[3]} 加到 \code{input[1]}。比较图
\ref{fig:10-7} 与图 \ref{fig:10-6} 的运算和操作数次序可以发现，这里实际上重排了
输入值列表，而不只是改变括号位置。有效次序变为

\[
  \{\code{input[0]},\code{input[8]},\code{input[1]},\code{input[9]},
  \ldots,\code{input[7]},\code{input[15]}\}.
\]

要使这种重排后的结果始终不变，运算符除了满足结合律，还必须满足交换律。

\begin{figure}[htbp]
  \centering
  \begin{minipage}{0.95\textwidth}
  \begin{lstlisting}[language=C++,numbers=left]
__global__ void ConvergentSumReductionKernel(float* input, float* output) {
    unsigned int i = threadIdx.x;
    for(unsigned int stride = blockDim.x; stride >= 1; stride /= 2) {
        if(threadIdx.x < stride) {
            input[i] += input[i + stride];
        }
        __syncthreads();
    }
    if(threadIdx.x == 0) {
        *output = input[0];
    }
}
  \end{lstlisting}
  \end{minipage}
  \caption{控制流分歧更少、执行资源利用效率更高的归约内核}
  \label{fig:10-8}
\end{figure}

图 \ref{fig:10-8} 实现了图 \ref{fig:10-7} 的工作分配。它与图 \ref{fig:10-5}
有几处细微但关键的区别。第一，所有者位置 \code{i} 设为 \code{threadIdx.x}，而不是
\code{2*threadIdx.x}（第 2 行），相邻线程的所有者位置因而也相邻。第二，步长从
\code{blockDim.x} 开始，每次减半，直到 1（第 3 行）。第三，每次迭代只有下标小于
步长的线程保持活跃（第 4 行），所以所有活跃线程的下标始终连续。

这些改动使第一轮不再把相邻元素相加，而是把相隔半个区段的元素相加；区段大小始终是
剩余活跃线程数的两倍。第一轮中各对元素都相距 \code{blockDim.x}，迭代结束后，所有
成对部分和都保存在输入数组前半部分。循环随后把步长除以 2，因此第二次迭代的步长为
\code{blockDim.x} 的一半，活跃线程所加的元素相隔四分之一个区段。

图 \ref{fig:10-8} 的循环中仍有 if 语句，每次迭代执行加法的线程数也与图
\ref{fig:10-5} 相同，为什么分歧会减少？答案在于执行加法的线程与不执行加法的线程
如何分布。以 256 个输入元素为例，第一次迭代所有线程都活跃，没有分歧。第二次迭代
中，线程 0--63 执行加法，线程 64--127 不执行。warp 由 \code{threadIdx.x} 连续的
32 个线程组成，因此 warp 0 和 warp 1 中的所有线程都执行加法，warp 2 和 warp 3
全部不执行；每个 warp 内的线程走同一条控制路径，所以没有控制流分歧。

不过，这个内核没有彻底消除 if 语句造成的分歧。对 256 元素的例子，从第 4 次迭代
开始，执行加法的线程少于 32；最后五次迭代分别只有 16、8、4、2 和 1 个线程执行加法，
因而仍有分歧。但存在分歧的循环迭代数已从原内核中的全部后续阶段减少为 5 次。执行资源
总消耗为

\[
  \left(\frac{N}{64}+\frac{N}{64}\cdot\frac12+\cdots+1+5\right)\cdot32.
\]

括号前半部分表示每次迭代都有一半 warp 完全退出，不再消耗资源，直到只剩一个完整
warp；最后一项 5 表示最后五次迭代中都只有一个活跃 warp，尽管其中仅有部分线程工作，
32 个线程对应的资源仍全部被占用。对 256 个输入元素，总消耗为
$(4+2+1+5)\times32=384$。活跃线程数没有改变，所以图 \ref{fig:10-8} 的效率为
$255/384=66\%$，几乎是图 \ref{fig:10-5} 的两倍。warp 会在执行资源有限的 SM 上
轮流调度，降低资源消耗也会缩短总执行时间。

\noindent\textbf{译者注：}底本此处把图 \ref{fig:10-5} 的资源消耗写成 736，但按
前文公式应为 $(4\times6+2+1)\times32=864$。本译稿沿用前文可复算的 864；384
约为 864 的 44\%，仍支持“显著减少”的结论。

图 \ref{fig:10-5} 与图 \ref{fig:10-8} 的代码差别很小，性能影响却可能很大。要有把握
做出这种调整，必须清楚理解设备的 SIMD 硬件如何执行线程。

\section{减少内存访问分歧}
\label{sec:reduction-memory-divergence}

图 \ref{fig:10-5} 的简单内核还有一个性能问题：内存访问分歧。第 5 章已经说明，每个
warp 内应尽量实现内存访问合并，也就是相邻线程访问全局内存中的相邻位置。但图
\ref{fig:10-6} 中，相邻线程并不访问相邻位置。每次迭代中，每个线程执行两次全局内存
读取和一次写入：第一次读所有者位置，第二次读相距 \code{stride} 的位置，最后写回
所有者位置。相邻线程拥有的位置并不相邻，所以访问无法充分合并；一次迭代中，一个
warp 集体访问的位置彼此相隔一个步长。

例如，第一次迭代中一个 warp 的线程执行第一次读取时，访问位置彼此相隔两个元素，
因此触发的全局内存事务数是合并访问的两倍，返回的数据也只有一半被线程使用。第二次
读取和写入也有同样问题。第二次迭代中，每隔一个线程就会退出，warp 集体访问的位置
相隔四个元素，于是需要四倍的全局内存事务，返回数据只有四分之一被使用。步长越大，
访问效率就越低。

这个过程一直持续到每个仍活跃的 warp 只剩一个活跃线程。此时一个 warp 才只发出一次
全局内存请求。因此，全局内存请求总数为

\[
  \left(\frac{N}{64}\cdot5\cdot2+
  \frac{N}{64}+\frac{N}{64}\cdot\frac12+
  \frac{N}{64}\cdot\frac14+\cdots+1\right)\cdot3.
\]

第一项对应前五次迭代：全部 $N/64$ 个 warp 中都至少有两个活跃线程，分散访问会触发
两个全局内存请求。其余项对应最后几次迭代：每个活跃 warp 只剩一个活跃线程，每次只
触发一个请求，而活跃 warp 数逐次减半。最后乘以 3，是因为每个活跃线程在一次迭代中
有两次读和一次写。对 256 个元素，图 \ref{fig:10-5} 的内核共发出
$(4\times5\times2+4+2+1)\times3=141$ 次全局内存请求。

图 \ref{fig:10-8} 的内核中，同一 warp 的相邻线程始终访问相邻的全局内存位置，所以
访问始终能够合并。每个 warp 的读写只触发最少数量的全局内存事务，并充分利用返回
数据。迭代推进时，整个 warp 逐批退出，非活跃 warp 不再访问全局内存；其余活跃 warp
继续发出合并访问，直到最后五次迭代只剩一个 warp。仅当请求的数据量小于单次全局内存
事务的大小时，才会出现资源利用不足。DRAM 带宽有限，因此图 \ref{fig:10-8} 内核的
执行时间很可能显著短于图 \ref{fig:10-5} 的简单内核。总之，前者既减少了控制流分歧，
又改善了内存访问合并，能更高效地利用执行资源和 DRAM 带宽。

\noindent\textbf{译者注：}底本本段两次写成“图 10.7 的内核”，但图 10.7 是工作
分配示意，实际内核位于图 10.8。本译稿将引用改为图 \ref{fig:10-8}。

\section{减少全局内存访问}
\label{sec:reduction-global-accesses}

还可以用共享内存进一步改进图 \ref{fig:10-8} 中分歧较少的内核。每次迭代都会把部分和
写到全局内存，下一次迭代中同一线程或其他线程又会把它们读回来。虽然部分读写可能由
末级缓存提供服务，但访问末级缓存仍受明显的延迟和带宽限制。共享内存比末级缓存和全局
内存延迟低、带宽高，因此把部分和保留在共享内存中，可以继续提高执行速度。图
\ref{fig:10-9} 展示了这种思路。

\begin{figure}[htbp]
  \centering
  \small
  \begin{tabular}{c@{\quad}l}
    \textbf{全局内存} &
      \colorbox{blue!25}{\strut\quad\quad\quad\quad\quad\quad\quad\quad\quad}\quad
      初次加载两个输入元素\\[0.55em]
    $\Downarrow$ & 每个线程先把两个输入相加\\[0.55em]
    \textbf{共享内存} &
      \colorbox{green!30}{\strut\quad\quad\quad\quad}\quad
      后续部分和始终在共享内存中读写\\[0.45em]
    & \colorbox{green!30}{\strut\quad\quad}\quad$\rightarrow$
      \colorbox{green!30}{\strut\quad}\quad$\rightarrow$
      \colorbox{green!30}{\strut\ }
  \end{tabular}
  \caption{使用共享内存减少全局内存访问}
  \label{fig:10-9}
\end{figure}

图 \ref{fig:10-10} 的内核实现了这项策略。每个线程先从全局内存加载两个原始元素并相加，
再把部分和写入共享内存（第 4 行）。第一次归约已经在循环外访问全局内存时完成，所以
for 循环从 \code{blockDim.x/2} 开始，而不是 \code{blockDim.x}。把
\code{__syncthreads()} 移到循环开头，是为了在初次写入共享内存与第一次循环迭代之间
同步。后续迭代只读写共享内存（第 8 行）。最后，线程 0 把总和写入 \code{output}，
与前几个内核保持相同行为（第 11--13 行）。

\begin{figure}[htbp]
  \centering
  \begin{minipage}{0.95\textwidth}
  \begin{lstlisting}[language=C++,numbers=left]
__global__ void SharedMemorySumReductionKernel(float* input, float* output) {
    __shared__ float input_s[BLOCK_DIM];
    unsigned int t = threadIdx.x;
    input_s[t] = input[t] + input[t + BLOCK_DIM];
    for(unsigned int stride = blockDim.x/2; stride >= 1; stride /= 2) {
        __syncthreads();
        if(threadIdx.x < stride) {
            input_s[t] += input_s[t + stride];
        }
    }
    if(threadIdx.x == 0) {
        *output = input_s[0];
    }
}
  \end{lstlisting}
  \end{minipage}
  \caption{用共享内存减少全局内存访问的归约内核}
  \label{fig:10-10}
\end{figure}

这样，全局内存访问只剩初次加载原输入以及最后写出结果。归约 $N$ 个元素时，全局内存
访问次数为 $N+1$。图 \ref{fig:10-10} 第 5 行的两组全局内存读取都能够合并，因此只
发出 $(N/32)+1$ 次全局内存请求。对 256 个元素，图 \ref{fig:10-8} 的内核需要 36 次
请求，共享内存内核则只需 $8+1=9$ 次，改善 4 倍。共享内存除了减少全局内存访问，还有
一个好处：输入数组不会被修改。如果程序其他部分还要使用原始数组值，这项性质很重要。

\noindent\textbf{译者注：}底本图 10.10 的函数签名只列出 \code{input} 参数，函数体
却写入未声明的 \code{output}；紧接着的正文又把最后一次写出误写成写到
\code{input[0]}。根据图 10.5、图 10.8 以及图 10.10 函数体的一致行为，本译稿在签名中
补入 \code{float* output}，并把正文改为“写出结果”。

\section{用 warp 级原语减少同步开销}
\label{sec:reduction-warp-primitives}

图 \ref{fig:10-10} 的内核中，当步长减到 32 及以下时，线程块里只有一个 warp 仍参与
归约树计算。CUDA 提供多种\textbf{warp 级原语（warp-level primitive）}，让同一
warp 的线程以跨 warp 线程无法使用的方式共享数据和同步。因此，当活跃线程数降至 32
及以下时，可以用这些原语减少共享内存访问和 \code{__syncthreads()} 的开销。本节先
介绍相关 CUDA warp 级原语，再用它们继续优化归约内核。

最简单的 warp 级原语是 \code{__syncwarp()}。它的作用类似 \code{__syncthreads()}，
但范围只有一个 warp：warp 中所有线程都到达调用位置之后，任何线程才能继续执行。
如果 warp 中的线程此前发生分歧，\code{__syncwarp()} 会保证它们重新汇合。

与本节归约优化尤其相关的另一类原语是\textbf{warp shuffle 函数（warp shuffle
function）}。同一 warp 中的线程可以用它们直接交换（shuffle）寄存器里的数据，无须
访问共享内存，也无须调用 \code{__syncthreads()}。其中，\code{__shfl_down_sync}
能让线程取得下标比自己大一个给定步长的线程所持数据。图 \ref{fig:10-11} 说明了它的
参数。

\begin{figure}[htbp]
  \centering
  \begin{minipage}{0.90\textwidth}
  \begin{lstlisting}[language=C++,numbers=none]
T __shfl_down_sync(unsigned mask, T var, unsigned int delta);
  \end{lstlisting}
  \begin{itemize}
    \item \code{mask} 指定执行 shuffle 指令时 warp 中哪些线程处于活跃状态；
    \item \code{var} 是发送线程交给接收线程的值；
    \item \code{delta} 是发送线程下标相对接收线程下标高出的步长；
    \item 返回值是接收线程从发送线程取得的值。对线程 $i$，若线程 $i$ 与
    $i+\code{delta}$ 都在 \code{mask} 中，则返回线程 $i+\code{delta}$ 的
    \code{var}。
  \end{itemize}
  \end{minipage}
  \caption{warp 级原语 \code{__shfl_down_sync}}
  \label{fig:10-11}
\end{figure}

借助 \code{__shfl_down_sync}，可以在不访问共享内存、不调用
\code{__syncthreads()} 的情况下完成 warp 内归约。图 \ref{fig:10-12} 展示了如何把
这种归约接到块内归约的最后阶段，从而减少屏障同步与共享内存访问。为了画面简洁，图中
一个线程块由 4 个 warp 构成，每个 warp 只有 4 个线程；实际配置中，1024 线程的块会
包含 32 个 warp，每个 warp 有 32 个线程。

\begin{figure}[htbp]
  \centering
  \small
  \begin{tabular}{c|c|c|c|l}
    warp 0 & warp 1 & warp 2 & warp 3 &\\\hline
    \multicolumn{4}{c|}{\colorbox{blue!25}{\strut\hspace{20em}}} & 从全局内存加载到共享内存\\[0.45em]
    \colorbox{green!30}{\strut\hspace{4em}} &
    \colorbox{green!30}{\strut\hspace{4em}} &
    \colorbox{green!30}{\strut\hspace{4em}} &
    \colorbox{green!30}{\strut\hspace{4em}} & 块内迭代并同步\\[0.45em]
    \colorbox{green!30}{\strut\hspace{4em}} &
    \colorbox{green!30}{\strut\hspace{4em}} &&& 继续在共享内存中归约\\[0.45em]
    \colorbox{orange!35}{\strut\hspace{4em}} &&&& 剩余 warp 改用寄存器 shuffle\\[0.35em]
    \colorbox{orange!35}{\strut\hspace{2em}} &&&& $\Downarrow$\\[-0.1em]
    \colorbox{orange!35}{\strut\hspace{0.6em}} &&&& 最终结果
  \end{tabular}
  \caption{在最后阶段使用 warp 内归约，以减少同步开销和共享内存访问}
  \label{fig:10-12}
\end{figure}

图 \ref{fig:10-12} 的前几次迭代与图 \ref{fig:10-9} 相同。每个线程从全局内存加载
两个元素，相加后把结果存入共享内存；各次迭代之间用 \code{__syncthreads()} 执行块级
屏障同步。每轮结束后，一半线程退出。不同之处是，一旦剩余线程只构成一个 warp，代码
便不再使用 \code{__syncthreads()} 和共享内存，而是把共享内存中的值载入寄存器，通过
在寄存器之间 shuffle 数据来执行归约树的剩余迭代。

为了实现图 \ref{fig:10-12}，先定义两个设备辅助函数，分别确定线程块内的 warp 下标，
以及线程在 warp 内的下标：

\begin{lstlisting}[language=C++,numbers=none]
__device__ unsigned int warpIdx() { return threadIdx.x/WARP_SIZE; }
\end{lstlisting}

这个函数用线程下标除以 32，使每组连续 32 个线程得到同一个 warp 下标。例如，线程
0--31 的 warp 下标都是 0，线程 32--63 的 warp 下标都是 1，依此类推。线程在 warp
内的下标也称\textbf{lane 下标（lane index）}，定义为：

\begin{lstlisting}[language=C++,numbers=none]
__device__ unsigned int laneIdx() { return threadIdx.x%WARP_SIZE; }
\end{lstlisting}

它计算线程下标对 warp 大小 32 的余数，所以同一 warp 的连续 32 个线程的 lane 下标
为 0--31。例如，块内线程 32 是 warp 1 的第一个线程，lane 下标为 $32\bmod32=0$；
线程 33 是第二个，lane 下标为 $33\bmod32=1$。

有了这两个辅助函数，就可以实现 warp 内归约。图 \ref{fig:10-13} 给出的设备函数为
\code{warp_reduce}，它通过寄存器 shuffle 完成归约。warp 中所有线程都应调用它，
每个线程通过参数 \code{val} 提供自己要贡献的值。线程先以 \code{val} 初始化局部
部分和 \code{partialSum}（第 2 行），再执行归约循环，让步长从 16 逐次减半到 1
（第 3 行）。起始步长为 16，是因为一个 warp 有 32 个线程、需要归约 32 个值；第一次
迭代中，前 16 个线程分别累加后 16 个线程的值。

\begin{figure}[htbp]
  \centering
  \begin{minipage}{0.92\textwidth}
  \begin{lstlisting}[language=C++,numbers=left]
__device__ __inline__ float warp_reduce(float val) {
    float partialSum = val;
    for(unsigned int stride = WARP_SIZE/2; stride > 0; stride /= 2) {
        partialSum += __shfl_down_sync(0xffffffff, partialSum, stride);
    }
    return partialSum;
}
  \end{lstlisting}
  \end{minipage}
  \caption{使用 warp 级原语执行 warp 内归约的设备函数}
  \label{fig:10-13}
\end{figure}

归约循环的每次迭代都调用 \code{__shfl_down_sync}，以较低开销在 warp 内交换数据并
同步（第 4 行）。第一个参数 \code{0xffffffff} 是所有位均为 1 的 32 位掩码，表示
warp 中全部线程都参与归约。第二个参数表示各发送线程传出寄存器中的
\code{partialSum}；第三个参数表示接收线程的下标比发送线程低 \code{stride}。函数
返回接收线程从发送线程得到的值。也就是说，对 warp 中的线程 $i$，它返回线程
$i+\code{stride}$ 的 \code{partialSum}，再把该值加到线程 $i$ 的部分和上。循环结束
后，warp 第一个线程的 \code{partialSum} 才是最终结果。函数把值返回给调用方（第 6
行），但只有第一个线程的返回值代表真实归约结果。

\begin{figure}[htbp]
  \centering
  \begin{minipage}{0.97\textwidth}
  \begin{lstlisting}[language=C++,numbers=left]
__global__ void WarpLevelSumReductionKernel(float* input, float* output) {
    __shared__ float input_s[BLOCK_DIM];
    unsigned int t = threadIdx.x;
    input_s[t] = input[t] + input[t + BLOCK_DIM];
    for(unsigned int stride = blockDim.x/2; stride >= WARP_SIZE;
        stride /= 2) {
        __syncthreads();
        if(threadIdx.x < stride) {
            input_s[t] += input_s[t + stride];
        }
    }
    __syncthreads();
    if(warpIdx() == 0) {
        float partialSum = input_s[t];
        partialSum = warp_reduce(partialSum);
        if(threadIdx.x == 0) {
            *output = partialSum;
        }
    }
}
  \end{lstlisting}
  \end{minipage}
  \caption{在最后阶段采用 warp 内归约，以减少同步开销和共享内存访问的归约内核}
  \label{fig:10-14}
\end{figure}

图 \ref{fig:10-14} 用刚才的设备函数优化归约内核。前半部分（第 1--12 行）与图
\ref{fig:10-10} 类似，不过第 5 行的循环在步长达到 warp 大小 32 时结束共享内存阶段。
从剩余 32 个部分和开始，只需块内第一个 warp 继续执行（第 13 行）。这个 warp 的每个
线程把共享内存中与自身下标对应的元素载入寄存器 \code{partialSum}（第 14 行），再共同
调用图 \ref{fig:10-13} 的 \code{warp_reduce}（第 15 行）。最终和落到 warp 的第一个
线程，也就是块内线程 0，由它写入全局内存中的输出（第 16--18 行）。这样，最后几个
步长阶段不再访问共享内存，也不再执行块级屏障同步。

\code{__shfl_down_sync} 以及其他 warp 级原语还接受一个可选参数，用来指定协作线程组
的宽度。默认宽度为 32，即整个 warp 是一个组；也可以指定更小的宽度，把 warp 划成
多个独立协作的小组。例如宽度为 8 时，一个 warp 被视为 4 个 8 线程小组，每组内部
协作执行原语，各组互不影响。

CUDA 还提供 \code{__shfl_sync}、\code{__shfl_up_sync} 和 \code{__shfl_xor_sync}
等 warp shuffle 函数，适合需要在同一 warp 的不同线程间交换数据的其他场景。其具体
行为见 CUDA C++ 编程指南。除 shuffle 之外，warp 级原语还包括：检查 warp 中哪些
线程满足某个条件的 warp 投票函数；检查哪些线程的某个参数值相同的 warp 匹配函数；
以及对整数执行 warp 内归约的 warp 归约函数。第 12 章将展示 warp 投票函数的用法。

\section{用两阶段 warp 内归约进一步减少同步开销}
\label{sec:reduction-two-stage}

图 \ref{fig:10-14} 的内核消除了归约树最后几次迭代中的共享内存访问和屏障同步，但最初
几轮仍要承担这些开销。还可以在初始阶段也采用 warp 内归约：线程块中的每个 warp 分别
归约一段输入，把部分和写入共享内存，再由一个 warp 对这些部分和执行 warp 内归约。

\begin{figure}[htbp]
  \centering
  \small
  \begin{tabular}{c|c|c|c|l}
    warp 0 & warp 1 & warp 2 & warp 3 &\\\hline
    \colorbox{orange!35}{\strut\hspace{4em}} &
    \colorbox{orange!35}{\strut\hspace{4em}} &
    \colorbox{orange!35}{\strut\hspace{4em}} &
    \colorbox{orange!35}{\strut\hspace{4em}} & 各 warp 在寄存器中独立归约\\[0.5em]
    \colorbox{orange!35}{\strut\hspace{2em}} &
    \colorbox{orange!35}{\strut\hspace{2em}} &
    \colorbox{orange!35}{\strut\hspace{2em}} &
    \colorbox{orange!35}{\strut\hspace{2em}} & $\Downarrow$\\[0.5em]
    \multicolumn{4}{c|}{\colorbox{green!30}{\strut\hspace{8em}}} & 各 warp 写入共享内存并块级同步\\[0.5em]
    \colorbox{orange!35}{\strut\hspace{4em}} &&&& 第一个 warp 归约全部部分和\\[0.35em]
    \colorbox{orange!35}{\strut\hspace{1em}} &&&& 最终结果
  \end{tabular}
  \caption{用两阶段 warp 内归约进一步减少同步开销和共享内存访问}
  \label{fig:10-15}
\end{figure}

图 \ref{fig:10-15} 展示了这种两阶段方案。最初，每个 warp 加载分配给自己归约的元素，
并执行一次 warp 内归约。第一阶段结束时，每个 warp 的第一个线程保存该 warp 所有输入
元素之和。这些线程把各自的和写入共享内存，然后调用 \code{__syncthreads()} 执行一次
块级屏障同步。\footnote{这里仍假设只用一个线程块归约最多 2048 个元素，所以所有 warp
可以通过共享内存交换数据。下一节将去掉这项假设。}随后，第一个 warp 读取这些部分和，
再执行一次 warp 内归约，得到总和。

\begin{figure}[htbp]
  \centering
  \begin{minipage}{0.97\textwidth}
  \begin{lstlisting}[language=C++,numbers=left]
__global__ void TwoStageWarpLevelSumReductionKernel(float* input, float* output) {

    unsigned int t = threadIdx.x;
    float partialSum = input[t] + input[t + BLOCK_DIM];
    partialSum = warp_reduce(partialSum);

    __shared__ float partialSums_s[BLOCK_DIM/WARP_SIZE];
    if(laneIdx() == 0) {
        partialSums_s[warpIdx()] = partialSum;
    }
    __syncthreads();

    if(warpIdx() == 0) {
        float partialSum = (t < BLOCK_DIM/WARP_SIZE) ? partialSums_s[t] : 0.0f;
        partialSum = warp_reduce(partialSum);
        if(threadIdx.x == 0) {
            *output = partialSum;
        }
    }
}
  \end{lstlisting}
  \end{minipage}
  \caption{用两阶段 warp 内归约进一步减少同步开销和共享内存访问的归约内核}
  \label{fig:10-16}
\end{figure}

图 \ref{fig:10-16} 的内核用两阶段 warp 内归约完成块内归约。每个线程先加载分配给
自己的两个输入元素，相加后把结果放在寄存器中（第 4 行），同一 warp 的线程再共同调用
\code{warp_reduce}（第 5 行）。此时块内所有 warp 都活跃，各自执行独立的归约。随后，
每个 warp 中 lane 下标为 0 的第一个线程保存了本 warp 的部分和，它把这个值写入共享
内存中与 warp 对应的位置（第 7--10 行）。一次屏障同步保证所有 warp 都完成写入
（第 11 行）。

部分和全部就绪后，代码的后半段与图 \ref{fig:10-14} 类似。只有第一个 warp 继续执行
（第 13 行）。\footnote{一个线程块最多有 1024 个线程，而一个 warp 有 32 个线程，
因此块内最多有 32 个 warp，第一阶段也最多产生 32 个部分和；一个 warp 足以完成第二
阶段。}第一个 warp 的线程分别从共享内存加载部分和；部分和不足 32 个时，其余 lane
以加法单位元 $0.0$ 参与（第 14 行）。这些线程执行 warp 内归约（第 15 行），最后由
第一个线程把总和写入全局内存（第 16--18 行）。

\noindent\textbf{译者注：}本节以单线程块归约最多 2048 个元素为主要示例，底本图
10.16 因而可能隐含 \code{blockDim.x} 与 \code{BLOCK_DIM} 都等于 1024；在这一前提下，
\code{partialSums_s} 恰有 32 个元素，原代码不会越界。不过，内核本身没有明示只能按
这一配置启动。若 \code{blockDim.x} 与 \code{BLOCK_DIM} 相等但小于 1024，第一个 warp
中的多余 lane 读取
\code{partialSums_s[t]} 时会越界。本译稿用加法单位元 $0.0$ 填充这些 lane，在不改变
1024 线程情形的同时，把适用范围扩展到 32--1024 之间、大小为整 warp 倍数且
\code{blockDim.x} 与 \code{BLOCK_DIM} 相等的线程块；这是兼容性扩展，并非对 1024
线程示例的纠错。

两阶段实现的主要优点是第一阶段也不再承担共享内存访问和屏障同步开销。只有从第一阶段
过渡到第二阶段时，才需要访问一次共享内存并执行一次屏障同步。不过，相比图
\ref{fig:10-14}，它也有缺点。图 \ref{fig:10-14} 的第一阶段中，每次迭代都会有一半
warp 完全退出，而且剩余 warp 内没有额外的跨 warp 分歧；图 \ref{fig:10-16} 的两阶段
实现则让所有 warp 在整个第一阶段都保持活跃，但每次 shuffle 并非每个线程都做有用
工作。因此，它用更多的 warp 内资源闲置换取更少的共享内存访问和屏障同步。

实践中，这项权衡通常值得。数据从全局内存载入后，归约内核其余部分的性能主要受延迟
限制。消除共享内存访问与屏障同步，可以减少内核中的长延迟指令。相比之下，消除分歧
warp 虽然有效，但对延迟受限内核的收益未必明显，因为省下的执行周期很可能会被停顿
周期取代。

\section{任意长度输入的归约}
\label{sec:reduction-arbitrary-length}

目前的所有内核都假设只启动一个线程块，主要原因是它们用 \code{__syncthreads()} 在
所有活跃线程之间执行屏障同步，而该函数只能同步同一个块中的线程。这就把当前硬件上的
并行度限制为 1024 个线程。大型输入数组可能有数百万乃至数十亿个元素，启动更多线程
可以继续加速归约。由于没有适合在不同块的线程之间执行屏障同步的方法，必须允许不同块
相互独立地执行。

\begin{figure}[htbp]
  \centering
  \small
  \begin{tabular}{c}
    \colorbox{blue!25}{\strut\hspace{29em}}\\[-0.1em]
    \textbf{输入数组}\\[0.4em]
    $\swarrow\hspace{9em}\downarrow\hspace{9em}\searrow$\\[-0.1em]
    \begin{tabular}{c@{\hspace{3em}}c@{\hspace{3em}}c}
      \fbox{区段 0} & \fbox{区段 1} & \fbox{区段 $k-1$}\\[0.35em]
      $\Downarrow$ & $\Downarrow$ & $\Downarrow$\\[-0.1em]
      块 0 的归约树 & 块 1 的归约树 & 块 $k-1$ 的归约树\\[0.35em]
      $\searrow$ & $\downarrow$ & $\swarrow$
    \end{tabular}\\[-0.1em]
    \texttt{atomic add}\quad$\Downarrow$\\[-0.1em]
    \fbox{最终结果}
  \end{tabular}
  \caption{使用原子操作的多块归约}
  \label{fig:10-17}
\end{figure}

图 \ref{fig:10-17} 展示了使用原子操作进行多块归约的思路。先把输入数组划分为多个
区段，每个区段的大小适合一个线程块处理；所有块各自独立执行归约树，再用原子加法把
各块结果累加到最终输出。

图 \ref{fig:10-18} 实现了多块归约，与图 \ref{fig:10-16} 相比只有少量改动。首先，
根据线程的块下标给 \code{segment} 赋不同值（第 3 行），以完成区段划分。每个区段的
大小为 \code{2*blockDim.x}，也就是每块处理 \code{2*blockDim.x} 个元素。把区段大小
乘以块的 \code{blockIdx.x}，即可得到该块所处理区段的起始位置。例如，每块有 1024
个线程时，区段大小为 $2\times1024=2048$，块 0、块 1、块 2 的起始位置分别是 0、
2048 和 4096。

\begin{figure}[htbp]
  \centering
  \begin{minipage}{0.98\textwidth}
  \begin{lstlisting}[language=C++,numbers=left]
__global__ void reduce_kernel(float* input, float* output, unsigned int N) {

    unsigned int segment = 2*blockDim.x*blockIdx.x;
    unsigned int i = segment + threadIdx.x;
    float partialSum = input[i] + input[i + BLOCK_DIM];
    partialSum = warp_reduce(partialSum);

    __shared__ float partialSums_s[BLOCK_DIM/WARP_SIZE];
    if(threadIdx.x%WARP_SIZE == 0) {
        partialSums_s[threadIdx.x/WARP_SIZE] = partialSum;
    }
    __syncthreads();

    if(threadIdx.x < WARP_SIZE) {
        float partialSum = (threadIdx.x < BLOCK_DIM/WARP_SIZE) ? partialSums_s[threadIdx.x] : 0.0f;
        partialSum = warp_reduce(partialSum);
        if(threadIdx.x == 0) {
            cuda::atomic_ref<float, cuda::thread_scope_device>
                output_ref(*output);
            output_ref.fetch_add(partialSum, cuda::memory_order_relaxed);
        }
    }
}
  \end{lstlisting}
  \end{minipage}
  \caption{使用原子操作的多块求和归约内核}
  \label{fig:10-18}
\end{figure}

确定每块的起始位置后，块内所有线程就可以把自己的区段当作完整输入来处理。把
\code{threadIdx.x} 加到块的区段起始位置，得到每个线程的所有者位置（第 4 行）。局部
变量 \code{i} 保存线程在全局输入数组中的所有者位置，取代前面内核里的 \code{t}；
第 5 行也相应改用 \code{i} 访问全局输入。

执行块内归约的其余代码保持不变，这里可以换入任意一种单块归约实现。最后一项差异是：
每个块的线程 0（首领线程）不再直接把部分和写成全局输出值，而是用原子加法把它累加到
输出中（第 18--20 行）。原子操作保证各线程块的首领线程不会在累加部分和时发生竞态。

用原子操作把各部分和累加为总和，有一个细微但重要的性质：线程块可以按任意相对次序
贡献结果，因为原子操作只保证互斥，并不规定顺序。因此，要用原子操作把归约推广到多个
区段，运算符必须同时满足交换律和结合律。

也可以不用原子操作。各线程块的首领线程可以按块下标把结果写入部分和数组，再启动一个
只有单块的新内核生成最终结果；还可以把部分和数组复制到主机内存，由 CPU 完成归约。
在原子操作不可用或代价过高的系统中，用单块网格或主机完成最后一步更合适。这两种方法
也让程序员更容易控制部分和的累加次序。

\noindent\textbf{译者注：}底本图 10.18 的参数 \code{N} 没有在函数体中使用。底本还以
\code{BLOCK_DIM} 表示预期块宽；以下讨论假定它等于 \code{blockDim.x}。在此前提下，
代码只能处理若干完整区段。每个区段包含 \code{2*BLOCK_DIM} 个元素，因此输入长度必须是
该区段大小的整数倍。这里的
“任意长度”是指输入不再受单线程块 2048 个元素的上限约束，而不是已经处理了末尾不完整
区段。底本练习 5 专门要求为非整段长度补上
边界处理，本译稿保留这项递进关系，没有替练习预先补写。图中第二阶段的部分和加载沿用
图 \ref{fig:10-16} 的单位元填充兼容性扩展，适用条件也与该图相同。此外，调用方必须先
把 \code{*output} 初始化为加法单位元 $0.0$；否则原子累加会叠加到原有值上。底本没有
明说这项前置条件。

\section{用线程粗化降低开销}
\label{sec:reduction-thread-coarsening}

目前的归约内核都试图用尽可能多的线程来最大化并行度：归约 $N$ 个元素时启动 $N/2$
个线程。每块 1024 个线程时，需要 $N/2048$ 个线程块。然而，处理器的执行资源有限，
硬件可能只能并行执行其中一部分块；多出的块会被串行调度，每当一个旧块完成后再执行
一个新块。

为了并行化归约，我们把工作分散到多个线程与线程块，实际上付出了很高代价。前几节
已经看到，归约树每进入下一阶段，更多线程都会闲置，硬件利用率随之降低；而且每个
启动的线程块都会经历这段低利用率阶段。如果各块确实并行运行，这项代价不可避免；
但如果硬件原本就要把它们串行化，不如由程序以更高效的方式自行串行化。第 6 章介绍过
\textbf{线程粒度粗化（thread granularity coarsening）}，简称线程粗化：把一部分工作
串行地交给更少线程，降低并行化开销。下面先把更多元素分给每个线程块，以展示线程粗化
后的并行归约，再说明它如何减轻硬件利用不足。

\begin{figure}[htbp]
  \centering
  \small
  \begin{tabular}{c}
    \textbf{32 个输入元素，每个线程负责 4 个}\\[0.25em]
    \colorbox{blue!25}{\strut\hspace{28em}}\\[-0.05em]
    $T_t:\ \code{input[t]},\ \code{input[t+8]},\
      \ \code{input[t+16]},\ \code{input[t+24]}\quad(0\le t<8)$\\[0.2em]
    $\Downarrow\quad$各线程独立执行三步累加$\quad\Downarrow$\\[0.25em]
    \colorbox{green!30}{\strut\hspace{12em}}\\[-0.05em]
    $\Downarrow$\quad{}块内归约树$\quad\Downarrow$\\[0.25em]
    \fbox{8 个线程部分和}$\rightarrow$\fbox{4 个}$\rightarrow$
    \fbox{2 个}$\rightarrow$\fbox{总和}
  \end{tabular}
  \caption{归约中的线程粗化}
  \label{fig:10-19}
\end{figure}

图 \ref{fig:10-19} 把线程粗化应用于图 \ref{fig:10-9} 的例子。图 \ref{fig:10-9}
中，每个块接收 16 个元素，也就是每个线程处理 2 个元素。每个线程从全局内存加载两个
元素时独立把它们相加，然后所有线程协作执行归约树。图 \ref{fig:10-19} 的线程块按因子
2 粗化，于是每块接收两倍元素，即 32 个元素，每个线程处理 4 个。线程先在加载时独立
累加自己的 4 个元素，再共同执行归约树。累加 4 个元素需要的三个步位于图的最前面，
此时所有线程都活跃并完成有效工作；各线程独立处理自己的元素，无须同步。剩余的归约树
步骤与粗化前相同。

\begin{figure}[htbp]
  \centering
  \begin{minipage}{0.97\textwidth}
  \begin{lstlisting}[language=C++,numbers=left]
__global__ void CoarsenedSumReductionKernel(float* input, float* output) {
    ...
    unsigned int segment = COARSE_FACTOR*2*blockDim.x*blockIdx.x;
    unsigned int i = segment + threadIdx.x;
    float partialSum = input[i];
    for(unsigned int c = 1; c < COARSE_FACTOR*2; ++c) {
        partialSum += input[i + c*BLOCK_DIM];
    }
    ...
}
  \end{lstlisting}
  \end{minipage}
  \caption{采用线程粗化的求和归约内核}
  \label{fig:10-20}
\end{figure}

图 \ref{fig:10-20} 展示了多块归约内核中实现线程粗化的关键代码；底本用省略号表示与
图 \ref{fig:10-18} 相同的块内归约和原子累加部分。两者有两项主要区别。第一，计算块
内区段起点时乘以 \code{COARSE_FACTOR}，反映每个区段扩大了该倍数（第 3 行）。第二，
不再只把线程负责的两个元素相加，而是用粗化循环按 \code{COARSE_FACTOR} 遍历并累加
这些元素（第 6--8 行）。整个粗化循环中所有线程都活跃，部分和累加到局部变量
\code{partialSum}；线程彼此独立，所以循环内不调用 \code{__syncthreads()}。

\noindent\textbf{译者注：}底本图 10.20 的函数定义漏写返回类型 \code{void}，本译稿
按本章其他 \code{__global__} 内核的声明方式补上。代码中的两处省略号属于底本原有的
差异展示，并非译稿删节。

\begin{figure}[htbp]
  \centering
  \scriptsize
  \begin{tabular}{c@{\hspace{4em}}c}
    \begin{tabular}{c}
      \textbf{两个原始块由硬件串行执行}\\[0.25em]
      \colorbox{blue!25}{\strut\hspace{11em}}\quad 全部活跃\\
      $\Downarrow$ 归约树，逐步闲置 $\Downarrow$\\
      \fbox{块 0 完成}\\[0.35em]
      \colorbox{blue!25}{\strut\hspace{11em}}\quad 全部活跃\\
      $\Downarrow$ 归约树，逐步闲置 $\Downarrow$\\
      \fbox{块 1 完成}\\[0.25em]
      (a) 共 8 步：2 步充分利用，6 步利用不足
    \end{tabular}
    &
    \begin{tabular}{c}
      \textbf{一个粗化块完成两个原始块的工作}\\[0.25em]
      \colorbox{blue!25}{\strut\hspace{16em}}\\
      $\Downarrow$ 三步独立累加，全部线程活跃 $\Downarrow$\\
      \colorbox{green!30}{\strut\hspace{8em}}\\
      $\Downarrow$ 三步归约树，逐步闲置 $\Downarrow$\\
      \fbox{粗化块完成}\\[0.25em]
      (b) 共 6 步：3 步充分利用，3 步利用不足
    \end{tabular}
  \end{tabular}
  \caption{比较采用与不采用线程粗化的并行归约}
  \label{fig:10-21}
\end{figure}

图 \ref{fig:10-21} 比较了两种执行方式：图 \ref{fig:10-21}(a) 是两个未粗化的原始
线程块被硬件串行执行；图 \ref{fig:10-21}(b) 是一个粗化块完成两个原始块的工作。
在图 \ref{fig:10-21}(a) 中，第一个块先执行一个步，每个线程把自己负责的两个元素相加。
此时所有线程都活跃，硬件得到充分利用。接下来的三个步执行归约树，每一步都有一半线程
退出，硬件利用不足；每一步还需要屏障同步和共享内存访问。第一个块完成后，硬件再调度
第二个块，在另一个数据区段重复同样步骤。两个块合计需要 8 个步，其中只有 2 个步充分
利用硬件，另外 6 个步利用不足并需要屏障同步和共享内存访问。

图 \ref{fig:10-21}(b) 用一个按因子 2 粗化的线程块处理相同数据。最初三个步中，每个
线程累加自己负责的 4 个元素；三个步里所有线程都活跃，硬件始终得到充分利用，也没有
屏障同步或共享内存访问。剩余三个步执行归约树，每一步有一半线程退出，硬件利用不足，
并需要同步和共享内存。总计只需 6 个步而不是 8 个；充分利用硬件的步从 2 个增至 3 个，
利用不足且需要同步与共享内存的步从 6 个减至 3 个。因此，线程粗化有效降低了硬件利用
不足、同步和共享内存访问带来的开销。

理论上，粗化因子可以远大于 2；但线程越粗，实际并行执行的工作越少。提高粗化因子会
减少硬件利用的数据并行度。如果因子大到启动的线程块少于硬件可以同时执行的块数，就
无法充分利用可用的并行执行资源。最佳粗化因子要保证有足够线程块占满硬件，通常取决于
输入总长度和具体设备特征。

\section{小结}
\label{sec:reduction-summary}

并行归约模式在许多数据处理应用中发挥关键作用。顺序代码虽然简单，要获得高性能的并行
执行却需要多项技术：重新分配线程下标以减少分歧；用共享内存减少全局内存访问；用 warp
级原语减少屏障同步和共享内存访问；用原子操作完成多块归约；以及线程粗化。归约还是
前缀和模式的重要基础；前缀和是并行化许多应用时的关键算法组件，也是下一章的主题。

实践中，程序员通常无须从头实现并行归约，因为 Thrust\cite{bell2012thrust,nvidia-thrust2025}
和 CUB\cite{nvidia-cub2025} 等 CUDA 库已经提供针对 GPU 高度优化的内置实现。不过，
并行归约仍是学习并行化与优化技术的极佳载体，而这些技术可以迁移到许多应用场景。

\phantomsection
\section*{练习}
\addcontentsline{toc}{section}{练习}

\begin{enumerate}[label=\arabic*.,leftmargin=2.7em]
  \item 对图 \ref{fig:10-5} 的简单归约内核，若元素数为 1024、warp 大小为 32，
  第 5 次迭代时块内有多少个 warp 会发生分歧？

  \item 对图 \ref{fig:10-8} 的改进归约内核，若元素数为 1024、warp 大小为 32，
  第 5 次迭代时有多少个 warp 会发生分歧？

  \item 修改图 \ref{fig:10-8} 的内核，采用下面所示的访问模式：

  \begin{center}
  \scriptsize
  \setlength{\tabcolsep}{2.2pt}
  \begin{tabular}{c|*{16}{c}}
    &0&1&2&3&4&5&6&7&8&9&10&11&12&13&14&15\\\hline
    线程所有者&&&&&&&&&$T_0$&$T_1$&$T_2$&$T_3$&$T_4$&$T_5$&$T_6$&$T_7$\\
    步 0&&&&&&&&&\colorbox{blue!25}{\strut$\Sigma$}&\colorbox{blue!25}{\strut$\Sigma$}&
      \colorbox{blue!25}{\strut$\Sigma$}&\colorbox{blue!25}{\strut$\Sigma$}&
      \colorbox{blue!25}{\strut$\Sigma$}&\colorbox{blue!25}{\strut$\Sigma$}&
      \colorbox{blue!25}{\strut$\Sigma$}&\colorbox{blue!25}{\strut$\Sigma$}\\
    步 1&&&&&&&&&&&&&\colorbox{blue!25}{\strut$\Sigma$}&
      \colorbox{blue!25}{\strut$\Sigma$}&\colorbox{blue!25}{\strut$\Sigma$}&
      \colorbox{blue!25}{\strut$\Sigma$}\\
    步 2&&&&&&&&&&&&&&&\colorbox{blue!25}{\strut$\Sigma$}&
      \colorbox{blue!25}{\strut$\Sigma$}\\
    步 3&&&&&&&&&&&&&&&&\colorbox{blue!25}{\strut$\Sigma$}
  \end{tabular}
  \\[0.35em]
  \begin{tabular}{c@{\quad}l}
    步 0 & $\code{input[8+k]}\leftarrow\code{input[8+k]}+\code{input[k]}$，$0\le k<8$\\
    步 1 & $\code{input[12+k]}\leftarrow\code{input[12+k]}+\code{input[8+k]}$，$0\le k<4$\\
    步 2 & $\code{input[14+k]}\leftarrow\code{input[14+k]}+\code{input[12+k]}$，$0\le k<2$\\
    步 3 & $\code{input[15]}\leftarrow\code{input[15]}+\code{input[14]}$
  \end{tabular}
  \end{center}

  \item 修改图 \ref{fig:10-20} 的内核，使它执行最大值归约而不是求和归约。

  \item 修改图 \ref{fig:10-20} 的内核，使它能处理不一定为
  \code{COARSE_FACTOR*2*blockDim.x} 整数倍的任意长度输入。为内核增加一个表示输入
  长度的参数 \code{N}。

  \item 假设要对下面的输入数组执行并行归约：
  \[
    \{6,2,7,4,5,8,3,1\}.
  \]
  分别采用以下内核时，画出每次迭代后数组内容的变化：
  \begin{enumerate}[label=\alph*.,leftmargin=2.2em]
    \item 图 \ref{fig:10-5} 中未经优化的内核；
    \item 图 \ref{fig:10-8} 中针对访问合并与分歧优化的内核。
  \end{enumerate}
\end{enumerate}

\begin{chapterbibliography}{9}
  \bibitem{bell2012thrust}
  N. Bell, J. Hoberock, Thrust: a productivity-oriented library for CUDA,
  in: \emph{GPU Computing Gems Jade Edition}, Elsevier, 2012, pp. 359--371.

  \bibitem{nvidia-thrust2025}
  NVIDIA, Thrust: the C++ parallel algorithms library, 2025.
  \url{https://nvidia.github.io/cccl/thrust/index.html}.

  \bibitem{nvidia-cub2025}
  NVIDIA, CUB, 2025.
  \url{https://nvidia.github.io/cccl/cub/index.html}.
\end{chapterbibliography}
